\documentclass[10pt,a4paper]{amsart}
\usepackage[margin=19mm]{geometry}
\usepackage{amsmath,amssymb,mathtools}
\usepackage[scr=rsfs,cal=euler]{mathalfa}
\usepackage{xfrac}
\usepackage[unicode,hidelinks,bookmarks=false]{hyperref}
\newcommand{\ie}{i.e.\ }
\newcommand{\cf}{cf.\ }
\newcommand{\resp}{respectively\ }
\newcommand{\Subsection}{\subsection}
\providecommand{\nobreakdash}{\nobreak\hbox{-}}
\setlength{\emergencystretch}{2em}
\multlinegap0pt
\usepackage{amssymb}
\newtheorem{proposition}{Proposition}
\newcommand{\C}{\mathbb{C}}
\title{Cubature from rational approximation}
\author{Gentian Zavalani}
\date{Source: arXiv:2607.17851v1 (CC BY 4.0)}
\begin{document}
\maketitle

\noindent\emph{Source and licence.} Cubature from rational approximation. Version arXiv:2607.17851v1 (CC BY 4.0), distributed by arXiv under the Creative Commons Attribution 4.0 International licence, http://creativecommons.org/licenses/by/4.0/, as recorded in the arXiv licence metadata for this exact version and shown on the abstract page https://arxiv.org/abs/2607.17851v1. Full licence notice: \texttt{licenses/arxiv-2607.17851-CC-BY-4.0.txt}.\par\smallskip

\noindent\textit{Editorial scope.} Counted unit: the article's proposition prop:dual (source0.tex lines 358--381). The article's complete original proof of it is delivered with it (source0.tex lines 383--408, one \texttt{proof} environment of 26 lines). No mathematics is written, repaired or re-derived: the statement and the proof are the article's own lines, in the article's own notation, and no retained line is substituted or re-flowed.  No further source-local context is needed: every cross-reference of every retained line resolves inside the two delivered spans.  Nothing was omitted from any retained span, so the package carries no omission record.  No retained line contains a citation, so the wrapper carries no bibliography and no bibliography entry.  No retained line contains a figure environment, an image-inclusion command or drawing code, so no image is delivered and the package declares no asset dependency.  Editorial formatting only, in addition: an AMS \texttt{amsart} preamble that restates the article's own declarations and macros used by the retained lines, namely the environments \texttt{proposition} on separate counters, together with the standard packages those lines need.  The retained environments print in this document as Proposition 1 in order of appearance, whereas the article prints the same items with the numbering of its own full document.  The complete native source of the article as distributed in the arXiv e-print for this exact version is bundled unchanged as \texttt{sources/205601/source0.tex}.
\begin{proposition}[Conditional rational exactness]
\label{prop:dual}
Let $z_1,\ldots,z_n$ be the distinct interior cubature nodes, where
$n$ is their number, and set
\[
  r_n(z)=\sum_{k=1}^n\frac{c_k}{z-z_k},
  \qquad D(z)=r_n(z)-H(z).
\]
Suppose that $D$ has distinct zeros $s_1,\ldots,s_n$ in
$\C\setminus\bar\Omega$ and that the Cauchy matrix
\[
  C_{kj}=\frac{1}{z_k-s_j}
\]
is nonsingular. For any $f$ analytic in a neighbourhood of
$\bar\Omega$, let
\[
  q(z)=\sum_{j=1}^n\frac{\alpha_j}{z-s_j}
\]
be the unique rational function of this form satisfying
$q(z_k)=f(z_k)$. Then
\[
  I(q)=I_n(q)=I_n(f).
\]
\end{proposition}

\begin{proof}
Since the poles of $q$ lie outside $\bar\Omega$, Fubini's theorem and
the Cauchy formula give
\[
\frac{1}{2i}\oint_\gamma q(\xi)H(\xi)\,d\xi
 =\frac{1}{2\pi i}\iint_\Omega
   \oint_\gamma \frac{q(\xi)}{\xi-z}\,d\xi\,dA(z)
 =\iint_\Omega q(z)\,dA(z).
\]
Residue calculus also gives
$I_n(q)=(2i)^{-1}\oint_\gamma q(z)r_n(z)\,dz$. Hence
\[
I(q)-I_n(q)
  =-\frac{1}{2i}\oint_\gamma q(z)D(z)\,dz.
\]
The exterior poles of $q$ are cancelled by the zeros of $D$, so $qD$
is analytic in the exterior of $\gamma$. Moreover,
\[
  q(z)=O(z^{-1}),\qquad D(z)=O(z^{-1}),
\]
and therefore $q(z)D(z)=O(z^{-2})$. Deforming $\gamma$ to a large
circle in the exterior shows that the large-contour contribution
vanishes. Thus
$I(q)=I_n(q)$, while the interpolation conditions give
$I_n(q)=I_n(f)$.
\end{proof}

\end{document}
